\chapter{Failed Experiments, Errors and Rectification History}

This chapter reproduces the continuous issue-resolution log
(\texttt{docs/DEVLOG.md}). Entries are taken from the actual commands, error
messages, logs and results of the development; no entry was invented. Times are
UTC on 11 November 2026.

{\footnotesize
\setlength{\tabcolsep}{3pt}
\begin{longtable}{@{}r L{1.5cm}L{3.0cm}L{3.0cm}L{3.0cm}L{2.4cm}L{2.0cm}@{}}
\caption{Issue-resolution log.}\label{tab:devlog}\\
\toprule
\# & Stage & Issue / decision & Root cause / evidence & Correction & Verification & Outcome \\
\midrule\endfirsthead
\toprule \# & Stage & Issue / decision & Root cause / evidence & Correction & Verification & Outcome \\ \midrule\endhead
\csname @@input\endcsname ../paper/generated/devlog_rows.tex 
\bottomrule
\end{longtable}}

\section{Most consequential issues}
\textbf{Clutter-biased centring (\#13).} Symptom: in the exploratory plots, the
``centred'' model input placed the person about 0.7\,m off-centre, together with
a large second cluster. Diagnosis: range histograms of all returns vs.\ moving
returns showed that static clusters at 2.5\,m and 4.5\,m carried about half of
the returns, while all moving returns were on the person. Alternatives
considered: DBSCAN clustering (per-window cost, parameter sensitivity) and
radius/statistical outlier removal (removes sparse limb returns). Correction:
centring on moving returns plus 1.2\,m gating. Verification: the
before/after figure and the ablation \emph{-- spatial gating}. All runs were
restarted so that every reported number uses the same pipeline.

\textbf{Validation dominated by one recording (\#16--\#17).} Symptom: validation
accuracy (79\%) far below test accuracy (96\%), a seed spread of 9 points and an
epoch-1 checkpoint for one model. Diagnosis: per-recording validation accuracy
showed a single 186-s squat recording (20\% of validation windows) predicted as
``jump''. Correction: at most 55 evenly spaced windows per validation recording.
Verification: validation and test accuracy now agree. The test set played no
part in the fix.

\textbf{Quantisation (\#11).} PyTorch dynamic INT8 quantisation lost 7 points of
accuracy and was \emph{slower} for these tiny layers. Weight-only per-channel
INT8 kept accuracy (within 0.05 points) and reduced storage by a factor of 3.6.
Only the INT8 storage size is claimed, not an INT8 speed-up.

\textbf{Leakage findings (\#24--\#26).} The audit did not reveal leakage in our
own protocols. It did reveal that the RadHAR benchmark cannot test session or
subject independence, and that the random-window protocol used in some prior work
leaks most of its test frames into training. This changed the emphasis of the
evaluation (MiliPoint subject-disjoint as the primary result) and became a
contribution of the paper.
